\chapter{不可靠邻居信息下的选择与价值修正控制}
\label{chap:work2}

第三章的邻居影响机制依据相对综合奖励选择参考对象，能够在相应可靠信息条件下改变合作演化过程。但奖励排序本身不检验消息是否真实。当部分发送者持续报告虚假的高奖励时，排序规则可能反复选择这些发送者，并通过动作一致性项影响自身 Q 表。接收者既可能选错参考对象，也可能因被抬高的奖励差而施加过大的修正。本章围绕这两个环节，设计基于局部观测的参考对象选择与连续修正控制。

本章的方法面向明确的消息失真模型，不把异常节点识别作为预先给定的能力。控制器只接收本地可获取的信息，输出参考对象和更新幅度；智能体仍通过自己的 Q 表在线学习并选择合作或背叛。模型定义、训练目标和实现条件在本章给出，性能及机制检验由第五章的实验承担。

\section{问题定义与方法定位}
\subsection{受扰消息模型}

环境中的真实动作、博弈收益与声誉按第三章的规则产生。设全体节点集合为 $\mathcal V$，每次独立运行开始时，从中无放回选取 $K=\lfloor\rho N\rfloor$ 个异常发送者，记其集合为 $\mathcal B$。集合在该次运行期间保持不变。节点 $j$ 对外发送 $\widehat m_j(t)=(\widehat a_j(t),\widehat R_j(t))$；正常节点报告真实值，异常节点按照指定规则修改消息。

基础攻击将异常发送者的报告设为
\begin{equation}
 \widehat a_j(t)=D,\qquad \widehat R_j(t)=1,\qquad j\in\mathcal B.
 \label{eq:w2-high-defect}
\end{equation}
报告奖励仍位于归一化奖励的合法区间，因此不能仅靠检测数值越界排除。异常节点的真实动作仍由自身策略产生，故式\eqref{eq:w2-high-defect}描述的是虚假广播，不是强制这些节点在真实博弈中始终背叛。后续对照还分别改变奖励报告或动作报告，以分析两类信息的作用。

攻击只影响 NI 使用的消息。自身真实奖励、个体 Q 表和用于状态构造的声誉汇总不被该攻击直接篡改。这个限定使本章能够识别附加价值更新对错误消息的敏感性，但也意味着结果不覆盖任意状态污染或全通道攻击。评价时可使用 $\mathcal B$ 统计异常来源，控制器执行时不能读取这一集合。

\subsection{消息失真的分离与状态对照}

为区分奖励排序、动作语义和持续性的作用，除式\eqref{eq:w2-high-defect}外，定义表\ref{tab:w2-message-types}中的消息条件。每个异常发送者在同一轮向所有接收者广播同一条消息；正常发送者始终报告真值。随机报告中的奖励与动作独立抽样，且不因该消息是否最终被参考而改变抽样过程。间歇性失真则保留相同的异常身份集合，只改变这些发送者当前是否执行伪造规则。

\begin{table}[!htbp]
\centering\small
\caption{异常发送者的消息规则；$a_j$与$R_j$为该轮真值}
\label{tab:w2-message-types}
\begin{tabular}{lll}
\toprule
条件 & 动作报告 & 奖励报告 \\
\midrule
持续高奖励背叛 & $D$ & $1$ \\
仅抬高奖励 & $a_j$ & $1$ \\
仅翻转动作 & 与$a_j$相反 & $R_j$ \\
随机消息 & $C,D$各以$1/2$抽取 & $U[0,1)$ \\
间歇高奖励背叛 & 开启时$D$，否则$a_j$ & 开启时$1$，否则$R_j$ \\
较小奖励抬高 & $D$ & $\min\{1,R_j+0.2\}$ \\
高奖励合作 & $C$ & $1$ \\
\bottomrule
\end{tabular}
\par\smallskip{\zihao{-5}\rmfamily 来源：作者，本文模型定义与实验配置，2026年。\par}
\end{table}

间歇条件以零起始步号$t$计，当$t\bmod1000<500$时开启，其余轮关闭；各异常发送者采用相同开启阶段。高奖励合作报告用于检验合作偏好与消息真实性的区别：即使其报告不真实，也可能推动合作。因此，本章将这些条件称为消息失真，不预设每一种失真都降低合作，更不把它们视为覆盖全部策略性攻击的集合。

默认模型还依赖未被污染的邻域声誉状态。为检验结论对这一观测条件的依赖，另定义自身上一动作状态$s_i(t)=a_i(t-1)$，初始状态设为$D$，执行本轮动作后令$s_i(t+1)=a_i(t)$。该对照保持两状态、两动作的Q表规模以及奖励和NI公式不变。它改变的是个体学习器的状态输入，不要求从受扰报告中重建可信声誉；相应结果需与默认状态分别报告。

\subsection{最高奖励选择的暴露范围}

全局异常比例与单个接收者接触异常消息的概率不同。设接收者有 $d$ 个互不相同的候选邻居，在固定抽取 $K$ 个异常发送者时，至少一个候选异常的概率为
\begin{equation}
 p_{\mathrm{expose}}=1-\frac{\binom{N-d}{K}}{\binom{N}{K}},
 \label{eq:w2-exposure}
\end{equation}
其中组合数按可选节点数理解，该概率未以接收者自身正常为条件。若近似视各候选独立地以概率 $\rho$ 异常，则得到 $1-(1-\rho)^d$。在本章$N=900$、$K=90$的固定抽取设置下，式\eqref{eq:w2-exposure}给出$d=4$与$d=12$时的概率分别约为0.3444与0.7199。该计算说明，较少的异常发送者也可能覆盖大量接收者；扩大感知范围在增加潜在参考信息的同时，也增加了接触异常消息的机会。

暴露不等于实际损害。异常消息是否被选中、是否具有正奖励优势以及是否进入价值修正，取决于完整的选择和门控过程。因此，实验需要分别记录候选暴露、参考来源和实际更新，而不能直接用 $\rho$ 代替每个环节的异常比例。

\subsection{与已有通信防御思路的关系}

评估消息可靠性并调节其影响已有相关研究。ADMAC利用局部观测和历史信息估计消息可靠程度，并通过可分解的动作偏好结构控制各消息的贡献\cite{yu2024admac}。本章不把这一总体思路作为首创，而将研究对象具体化为社会困境中的在线价值学习：控制器影响每轮 Q 值的附加更新，个体动作价值仍在部署环境中持续变化。

这种设置带来两个需要单独检验的问题。首先，一条消息当轮改变的价值较小，仍可能通过多轮重复使用影响学习轨迹；其次，提高群体福利的选择规则可能只是偏好合作消息，并未准确识别真假。因此，本章用完整学习轨迹训练控制器，并设置简单过滤、固定弱影响和合作消息优先等对照，以判断复杂选择是否提供了额外作用。

\section{接收端局部特征}
\subsection{自身经验与消息集合}

接收者 $i$ 在完成本轮标准 TD 更新后，获得自身真实动作 $a_i(t)$、综合奖励 $R_i(t)$、当前 Q 表，以及候选集合 $\Omega'_i$ 内的报告消息。它还维护两个自身奖励预测值 $\mu_i^C(t)$ 与 $\mu_i^D(t)$，分别反映之前执行两种动作时得到的局部反馈。二者初始化为0.5。

预测量在本轮 NI 完成之后才更新，仅修改实际执行动作对应的一项：
\begin{equation}
 \mu_i^{a_i(t)}(t+1)=0.95\mu_i^{a_i(t)}(t)+0.05R_i(t),
 \label{eq:w2-predictor}
\end{equation}
另一项保持不变。由此，当前消息评分使用的是截至上一轮的预测量。更新仅依赖自身真实经验，不使用邻居真实收益或异常标签。当某种动作较少被执行时，其预测量可能陈旧，因此下文将它作为经验参照，而非经过校准的可信度估计。

\subsection{五维候选特征}

对每条候选消息构造向量 $x_{ij}(t)\in\mathbb R^5$，其分量定义为
\begin{equation}
 \begin{aligned}
 x_{ij}^{(1)}(t)&=\widehat R_j(t)-R_i(t),\\
 x_{ij}^{(2)}(t)&=2\mathbb I\{\widehat a_j(t)=a_i(t)\}-1,\\
 x_{ij}^{(3)}(t)&=\left|\widehat R_j(t)-\operatorname{median}_{k\in\Omega'_i}\widehat R_k(t)\right|,\\
 x_{ij}^{(4)}(t)&=\left|\widehat R_j(t)-\mu_i^{\widehat a_j(t)}(t)\right|,\\
 x_{ij}^{(5)}(t)&=\mathbb I\{\widehat a_j(t)=C\}.
 \end{aligned}
 \label{eq:w2-features}
\end{equation}
奖励优势用于衡量报告表现与自身表现的差别；动作一致性描述双方行为关系；中位数偏差提供同轮消息之间的横向比较；经验残差则比较报告值与自身同动作历史反馈。在十二个候选下，中位数取排序后中间两个奖励值的平均。

第五个分量显式表示候选报告的合作语义。只使用动作一致性时，线性评分器对“与合作接收者一致”和“与背叛接收者一致”采用相同权重，不能普遍表示固定合作优先规则。加入该分量后，取 $w=(1,0,0,0,2)$，便可使任何合作报告优先于背叛报告，并在同类动作中按奖励排序，因为奖励值位于 $[0,1]$。这一设置使候选模型包含简单合作优先作为可表达的特例；它提供公平比较的基础，不预先保证训练会超过该规则。附录\ref{app:development-history}保留未含该分量的首轮开发结果，说明两种表示的来源。

这些特征提供的是选择线索，不能直接等同于真假判定。真实的高收益邻居也可能偏离多数消息；不同空间位置的局部群组构成不同，邻居真实奖励也可能偏离自身预测。控制器需要根据长期任务表现学习如何组合这些线索，而不能预先把某个特征超过阈值定义为必然欺诈。

\section{参考对象选择与门控修正}
\subsection{共享评分器}

使用共享线性评分器
\begin{equation}
 z_{ij}(t)=w^\mathsf T x_{ij}(t),\qquad
 j_i^*(t)\in\operatorname*{arg\,max}_{j\in\Omega'_i}z_{ij}(t),
 \label{eq:w2-score}
\end{equation}
其中 $w\in\mathbb R^5$。所有接收者使用同一组评分参数，但输入来自各自局部环境。精确并列时，使用固定随机流中的参考对象抽样值均匀选择。评分器不添加对所有候选相同的偏置，因为该偏置不会改变排序；评分权重整体乘以正数也不改变参考对象，因此参数存储个数不等于所有方向都具有独立的可辨识作用。

与最高奖励规则相比，式\eqref{eq:w2-score}允许报告奖励较高但偏离邻域或自身经验的消息获得较低评分。它仍只选择一个参考对象，不聚合全部邻居 Q 表，也不生成合作或背叛动作。

\subsection{连续影响强度}

在当前状态下，定义自身两种动作价值接近程度的代理量
\begin{equation}
 u_i(t)=\frac{1}{1+|Q_i(s_i(t),C)-Q_i(s_i(t),D)|}.
 \label{eq:w2-q-proxy}
\end{equation}
Q 值取自本轮标准 TD 更新之后、NI 修正之前。$u_i(t)$ 越大，表示两种动作的当前估计越接近，但它不是已校准的不确定性概率。

连续门控定义为
\begin{equation}
 g_i(t)=\sigma\left(b+v^\mathsf T x_{i,j_i^*}(t)+v_u u_i(t)\right),
 \qquad \sigma(z)=\frac{1}{1+\exp(-z)}.
 \label{eq:w2-gate}
\end{equation}
其中 $v\in\mathbb R^5$，$v_u$ 和 $b$ 为标量。评分器与门控合计存储12个共享参数。门控表示附加更新的连续缩放程度，不是以某个概率使用消息。模型可以保留参考对象但降低其影响，这与直接丢弃该消息的操作不同。

\subsection{局部归一化与价值修正}

为使控制器执行只依赖局部输入，本章采用接收者自身的最大消息奖励差作为归一化依据：
\begin{equation}
 D_i(t)=\max_{j\in\Omega'_i}|\widehat R_j(t)-R_i(t)|+0.01.
 \label{eq:w2-local-normalizer}
\end{equation}
这与第三章的全局归一化不同，因此实验中保留全局 NI 与局部 NI 两种基线，单独评价这一变化。对选定对象，定义正奖励优势 $A_i(t)=\max\{0,\widehat R_{j_i^*}(t)-R_i(t)\}$，附加更新为
\begin{equation}
 \Delta Q_i(t)=\kappa_{\max}g_i(t)\frac{A_i(t)}{D_i(t)}
 \left[2\mathbb I\{a_i(t)=\widehat a_{j_i^*}(t)\}-1\right].
 \label{eq:w2-correction}
\end{equation}
该改变量加到当前状态下已执行动作的 Q 表项，标准 TD 项保持不变。开发训练取 $\kappa_{\max}=0.5$。若选定对象的报告奖励不高于自身，附加更新为0，即使它的相对评分最高，也不会自动获得强化权限。

由于 $A_i(t)\leq D_i(t)$ 且 $0\leq g_i(t)\leq1$，有
\begin{equation}
 |\Delta Q_i(t)|\leq\kappa_{\max}g_i(t)\leq\kappa_{\max}.
 \label{eq:w2-single-step-bound}
\end{equation}
这个界约束单次附加更新的幅度。若门控恒为0或 $\kappa_{\max}=0$，在相同初始化与动作随机流下，动作和 Q 表轨迹退化为独立 Q-learning。另一方面，该界不保证选中的邻居是真实的，也不约束两种不同策略在长期交互中产生的全部轨迹差异，因而不能被解释为一般的拜占庭鲁棒保证。

\subsection{消息动作对价值偏好的作用方向}

选择与门控具有不同的作用范围，这一点可以直接从更新式看出。令 $d_i=Q_i(s_i,C)-Q_i(s_i,D)$ 为当前状态下的合作价值优势，并记 $h_i=\kappa_{\max}g_iA_i/D_i\geq0$。仅考虑本轮附加更新对 $d_i$ 的改变量，有
\begin{equation}
 \Delta d_i^{\mathrm{NI}}=
 \begin{cases}
 +h_i, & \widehat a_{j_i^*}=C,\\
 -h_i, & \widehat a_{j_i^*}=D.
 \end{cases}
 \label{eq:w2-preference-direction}
\end{equation}
例如，参考对象报告合作时，接收者若执行合作，就增加合作Q值；若执行背叛，就降低背叛Q值，两者均增加合作价值优势。报告背叛时同理可得相反方向。该关系对接收者当轮执行哪一种动作都成立，但仅描述固定当前状态下的NI项，不包含标准TD更新和之后环境变化的影响。

因此，门控只能缩放既定参考对象引入的方向，不能把背叛报告产生的附加修正反向变成合作强化。若异常高奖励背叛报告长期占据最高奖励位置，仅调节幅度最多能够直接抑制这一方向性作用；更换参考对象才能改变该项的方向。另一方面，合作消息优先本身就能形成明确的合作偏置，所以较高合作率并不能单独证明模型识别了真实消息。式\eqref{eq:w2-preference-direction}为选择、门控和合作优先对照提供了共同的机制解释。

\subsection{有界奖励下的价值范围}

单次附加更新的幅度界还可以用于限定个体Q值的范围。设综合奖励始终位于$[0,1]$，TD学习率为$0<\eta\leq1$，折扣因子为$0\leq\gamma<1$，每轮NI满足式\eqref{eq:w2-single-step-bound}。令$q_{\min}^{0}$和$q_{\max}^{0}$分别为全部初始Q表项的最小值与最大值，定义
\begin{equation}
 \begin{aligned}
 L_Q&=\min\left\{q_{\min}^{0},-\frac{\kappa_{\max}}{\eta(1-\gamma)}\right\},\\
 U_Q&=\max\left\{q_{\max}^{0},\frac{\eta+\kappa_{\max}}{\eta(1-\gamma)}\right\}.
 \end{aligned}
 \label{eq:w2-q-range}
\end{equation}
则在上述更新规则下，每个时刻的全部Q表项均位于$[L_Q,U_Q]$。证明可以直接用归纳法完成。假设更新前全部表项位于该区间，考虑本轮被更新的一项$q$。TD之后再加NI的结果$q^+$满足
\begin{equation}
 \begin{aligned}
 q^+&\geq(1-\eta)L_Q+\eta\gamma L_Q-\kappa_{\max}\geq L_Q,\\
 q^+&\leq(1-\eta)U_Q+\eta(1+\gamma U_Q)+\kappa_{\max}\leq U_Q.
 \end{aligned}
 \label{eq:w2-q-range-proof}
\end{equation}
未被更新的表项仍在原区间内，初始时刻也由定义满足该条件，故结论成立。该推导不要求邻居策略固定，也不要求某个参考对象必然真实，仅依赖奖励范围和单次附加更新界。

以本章设置为例，取$\eta=0.8$、$\gamma=0.9$、$\kappa_{\max}=0.5$并采用近零初始化，可得保守范围$[-6.25,16.25]$。该结果说明有界NI不会使数学更新中的Q值无限增长，但不能证明Q值收敛、合作率提高或受扰轨迹接近无扰轨迹。实际观察到的范围可以更小；本节不将这个保守界作为防御性能指标。

\subsection{附加修正与有效奖励的关系}

附加项作用于Q表，也可以从更新的代数形式理解。记本轮实际NI增量为$\Delta_i^{\mathrm{NI}}(t)$，标准TD目标中的后继最大值为$Q_i^{\mathrm{next}}$，则同一被更新表项满足
\begin{equation}
 Q_i^+=(1-\eta)Q_i+\eta\left[R_i(t)+\frac{\Delta_i^{\mathrm{NI}}(t)}{\eta}
                  +\gamma Q_i^{\mathrm{next}}\right].
 \label{eq:w2-effective-reward}
\end{equation}
因此，从本轮最终更新结果看，NI相当于加入有效奖励增量$\Delta_i^{\mathrm{NI}}/\eta$。这里的有效奖励只是解释更新的记号，不替代物理博弈收益或实验福利。该增量依赖当前消息、参考对象以及TD更新后的价值差距，并非一个预先给定的平稳奖励函数。

这一视角说明，保留标准TD步骤不代表附加社会信息只会加快原有学习过程。Ng等在马尔可夫决策过程中讨论了保持最优策略的奖励变换，其中基于势函数的特定差分形式提供了相应条件\cite{ng1999policy}。本文没有将NI构造成这一形式，也不据此主张原个体最优策略保持不变。研究目标是评价消息引导的价值修正对群体合作的作用；第四章的方向分析与第五章的合作偏好对照，正是为区分这种行为引导与单纯提高信息准确性。

\section{基于完整学习轨迹的参数训练}
\subsection{目标函数与数据划分}

令 $\theta=(w,v,v_u,b)$ 为共享参数。每次评价一个候选参数时，重新初始化全部个体 Q 表和声誉，在指定博弈与消息条件下执行完整的在线学习过程。控制器参数在单次轨迹内保持固定，而 Q 表持续变化。训练评价对象因此是参数对整个合作形成过程的作用。

设训练环境集合为 $\mathcal E$，每个环境 $e$ 的协同因子为 $r_e>1$，训练时长为 $T$。开发训练目标定义为
\begin{equation}
 J(\theta)=\frac{1}{|\mathcal E|}\sum_{e\in\mathcal E}
 \mathbb E_{\xi}\left[\frac{1}{T}\sum_{t=0}^{T-1}
 \frac{\overline P_e(t;\theta,\xi)}{5(r_e-1)}\right],
 \label{eq:w2-objective}
\end{equation}
其中 $\xi$ 表示初始化与交互中的随机性。环境采用等权汇总，包含 $r\in\{4.4,4.8\}$ 与干净、10\%持续虚假高奖励背叛消息的四种组合。

由式\eqref{eq:foundation-welfare-identity}，这里的标准化真实福利严格等于同轮合作率。因此，式\eqref{eq:w2-objective}实际上优化四种环境下的全程平均合作水平，不能将其表述为与合作率无关。使用全程平均也意味着学习早期和后期都进入目标；最终评价另行报告末段指标，观察速度与后期水平是否一致。

训练环境种子取自110至199，模型比较使用独立验证种子，最终测试种子1000至1029在方法冻结前保留。优化初始化种子与环境种子分别记录。相同环境种子对应相同初始 Q 表、异常节点和逐轮动作随机数，使候选之间的比较尽量减少不相关随机波动；不同代可更换训练环境种子，避免只适应单一初始配置。

\subsection{交叉熵搜索过程}

按完整交互回报搜索参数，是强化学习任务中已有的一类优化途径。例如，Salimans等研究了以进化策略直接搜索策略参数的方法\cite{salimans2017evolution}。本文采用其中不依赖轨迹导数的黑盒评价思路，但优化对象是附加信息控制器，分布更新采用交叉熵精英统计；并未采用该文的参数扰动梯度估计器，也不沿用其任务上的速度结论。

由于参考对象包含离散最大值选择，且每次评价涉及较长的群体演化轨迹，本章采用高斯候选分布进行黑盒参数搜索。具体过程使用交叉熵方法的精英更新思路\cite{deboer2005tutorial}：在当前参数分布中采样候选，运行仿真计算目标值，再用得分较高的候选更新下一代分布。该过程没有对仿真轨迹进行反向传播，不能与可微分端到端训练混称。

两轮开发训练均采用24个候选、12代搜索，每个候选在四个环境单元各运行10000步。每代从同一个环境种子产生配对比较，取前20\%候选，即5个候选为精英。设候选分布的均值和标准差为 $m_g,s_g$，精英的逐维均值和标准差为 $\bar\theta_g,\widehat s_g$，则
\begin{equation}
 m_{g+1}=0.2m_g+0.8\bar\theta_g,\qquad
 s_{g+1}=\max\{0.15,0.2s_g+0.8\widehat s_g\},
 \label{eq:w2-cem}
\end{equation}
其中最大值逐维计算，候选参数限制在 $[-8,8]$，每代包含当前均值作为候选。当前五维特征模型的初始评分权重为 $(1,0,0,0,2)$，初始门控斜率为0、偏置为 $\log(2/3)$，因此初始门控为0.4，等效固定影响强度为0.2。该强度由独立开发批次中对简单合作优先规则的比较确定。评分部分初始标准差为1，门控部分为2。首轮四维特征模型另用 $(1,0,-1,-0.5)$ 初始化评分，门控参数为0；其结果只归属于对应模型，不能视为当前模型的实验。上述数值均为新增研究的开发配置，而非已发表工作的参数。

每轮采用三个独立优化初始化，分别执行完整搜索，各产生3456次候选仿真。每次优化保存每代全部候选、参数、环境种子及逐次评价结果，并输出最终分布均值与末代最佳候选。由于不同代使用的环境种子不同，不直接用跨代最高训练分数挑选唯一参数。候选仍需在共同的独立验证环境中比较，检验短时训练对较长评价时长的迁移。

\section{实现过程与验证设计}
\subsection{同轮数据与局部执行}

一轮完整执行先生成所有节点的动作，再计算同轮真实博弈收益，更新声誉并进行标准 TD 学习；随后根据收到的消息构造候选特征、选择参考对象、计算门控并修正当前 Q 表项；最后更新自身经验预测量。动作、收益和报告中的时间索引保持一致。全体节点共享控制器参数，但不共享 Q 表或自身经验预测量。

\begin{figure}[!htbp]
\centering
% Vector diagram of one receiver's NI update; all arrows follow implemented timing.
\begin{tikzpicture}[font=\small,>=stealth,
 box/.style={draw,rounded corners=2pt,align=center,minimum height=1.0cm,text width=4.7cm},
 wide/.style={box,text width=8.0cm}]
\node[box] (own) at (-3.0,0) {自身动作与真实奖励\\经验预测及TD更新后的Q表};
\node[box] (msg) at (3.0,0) {候选邻居报告\\$\{(\widehat a_j,\widehat R_j):j\in\Omega'_i\}$};
\node[wide] (feature) at (0,-1.8) {构造五维消息特征 $x_{ij}$\\计算局部归一化量 $D_i$ 与Q值接近量 $u_i$};
\node[wide] (score) at (0,-3.4) {共享评分：$z_{ij}=w^{\mathsf T}x_{ij}$\\选取参考对象 $j_i^*$，精确并列时均匀选择};
\node[wide] (gate) at (0,-5.0) {连续门控：$g_i=\sigma(b+v^{\mathsf T}x_{i,j_i^*}+v_u u_i)$\\正奖励优势决定是否产生附加修正};
\node[wide] (update) at (0,-6.6) {将带动作一致性符号的 $\Delta Q_i$ 加入当前表项\\更新自身经验预测，进入下一轮在线学习};
\draw[->] (own.south) -- (feature.north west);
\draw[->] (msg.south) -- (feature.north east);
\draw[->] (feature) -- (score);
\draw[->] (score) -- (gate);
\draw[->] (gate) -- (update);
\end{tikzpicture}

\caption{单个接收者在本轮标准TD更新之后的局部信息处理流程。共享参数由训练阶段获得，执行时固定；自身经验预测使用上一轮保留值，并在本轮NI完成后更新。异常身份和邻居未篡改真值不进入该流程。}
\label{fig:w2-local-flow}
\par\smallskip{\zihao{-5}\rmfamily 来源：作者，本文模型定义与实验配置，2026年。\par}
\end{figure}

实现中将参考选择封装为只接受局部可观测输入的计算过程，异常掩码不作为其输入。异常身份由仿真层用于生成受扰报告和计算诊断指标。每个节点在每轮固定抽取探索、随机动作、贪心平局及参考对象平局所需的随机量，不因算法分支不同而改变后续动作随机流。

随机初始化、异常集合、动作抽样与随机消息分别使用独立随机子流。加入随机消息不会消耗动作选择所用的随机数，因此在相同状态定义下，独立Q-learning的轨迹不随消息失真类型变化。对间歇性失真，异常来源选择比例仍按发送者身份统计，在关闭阶段选中该发送者并不表示选中了被篡改的消息。这个统计口径与控制器实际接收到的内容需要区分。

令候选数为 $d$。线性评分和门控计算随 $d$ 线性增长，中位数采用排序实现，单个接收者的相应代价为 $O(d\log d)$。每个节点额外维护两个自身经验预测值，控制器共享参数数量固定。实际开销还包括候选消息的构造与排序，不能只用12个参数这一数量宣称其运行开销可以忽略；第五章通过同配置测量报告实际开销。

\subsection{必要对照及机制归因}

最高奖励选择、小固定强度和上尾过滤分别提供排序、影响减弱和简单鲁棒统计的参照。另设置合作消息优先规则：只要存在报告合作的候选，就从中选择最高报告奖励；否则仍按最高报告奖励选择。该规则使用了消息中的动作语义，用于检查训练所得的高合作水平是否可以由一个简单合作偏置解释。

另设置只允许合作方向附加更新的规则：参考对象仍按合作优先选择，但当其报告背叛时将门控置0。由式\eqref{eq:w2-preference-direction}，这一规则的单步NI合作价值改变量非负。这个性质不能保证整个系统的合作率单调提高，因为标准TD更新和状态变化仍然存在；它用于检验禁止背叛方向修正是否已经足以解释性能。

模型内部通过两类重新训练的消融区分选择与门控：仅学习门控时固定最高奖励排序；仅学习选择时将门控固定为0.4，对应已选固定强度0.2。完整模型与两类消融使用一致的环境和候选评价预算。固定强度对照还在开发集校准平均实际修正幅度，以检验单纯减少更新量能否解释差异；校准的汇总容差与逐条件限制在第五章分别报告。

实现检查覆盖局部特征与手工计算的一致性、精确并列选择、零门控退化、单步更新幅度界及同种子重现。训练表现则需要独立验证；通过实现检查只说明指定机制按定义运行，不证明它一定优于其他方法。若简单规则取得相同或更好的效果，应保留该结果并据此讨论方法的必要性。

\section{本章小结}

本章把不可靠邻居信息的使用分为参考对象选择和价值修正控制两个环节。控制器结合邻居报告的相对奖励、动作关系、横向偏差及自身经验进行评分，并通过连续门控调节附加 Q 值更新。训练对象为共享控制器，部署时个体 Q 表仍从头在线学习。方法的有效性取决于它在独立环境中是否优于固定弱影响、简单过滤和合作消息优先等对照；其结论也限定于明确定义的消息攻击通道，不能扩展为一般可信度识别或任意攻击下的理论鲁棒性。
